\section{Introduction}

Agents built on language models now plan, call tools, and delegate to other agents. The dominant
way we measure them is task success: did the ticket get refunded, did the migration finish, did the
buy order go through \citep{yao2024taubench,liu2023agentbench,zhou2023webarena}. Success rates are
useful but they collapse three very different competencies into one bit. An agent that never
encounters a fault, an agent that encounters a fault and silently routes around it, and an agent
that hits a fault, notices, and repairs it are all scored as ``1''. An agent that fails because a
tool returned a stale read and an agent that fails because it misread the objective are both
scored as ``0''.

The consequence is that we cannot answer the question practitioners actually ask when an agent
breaks in production: \emph{would it have noticed?} A reliability claim about an agent is a claim
about a closed loop---detect, diagnose, repair, verify---and current benchmarks only score the last
step's shadow, whether the final answer happened to be right. Worse, the loop can look closed when
it is not: our agents assert success in their own trajectory \emph{far} more often than they
achieve it (Sec.~\ref{sec:results}), so a judge that reads the transcript will over-credit them.

This paper asks a narrow, measurable question:

\begin{quote}\itshape
Can agent failures be detected early and recovered from reliably, and which failure signals predict
successful recovery in multi-step tool-using and multi-agent workflows?
\end{quote}

\noindent We answer it with \arb{}, a benchmark that extends failure \emph{counting} into failure
\emph{handling} under control. Our contributions:

\begin{enumerate}[leftmargin=1.4em]
\item \textbf{A signal-defined taxonomy} of seven classes---misuse (invalid arguments), stale state,
contradictory evidence, planning error, communication failure, objective drift, and unrecoverable
external error---each with a characteristic observation signature, a canonical remediation operator,
and an explicit recoverability status (Table~\ref{tab:taxonomy}). The taxonomy is defined by what
you can \emph{see} and what you would \emph{do}, not by which module we believe is at fault, which
keeps labels independent of any particular agent architecture.

\item \textbf{Four deterministic tool environments} with verifiable world state (refund desk,
dual-store migration, inventory accounting with an auditable invariant, three-role procurement with
a message bus). Outcomes are read from state, never from the agent's prose. Faults are injected on a
schedule aligned to the step that calls the targeted tool, and every trajectory stores expected
state, observed state, the mismatch, and the resulting harm.

\item \textbf{An inert-fault screen.} We run every (task, fault) pair with no recovery first.
\InertPct\% of injections leave the task outcome unchanged. These are excluded from detection and
recovery denominators, so a model is never rewarded for ``detecting'' a fault that did not matter
and recall is not inflated by faults no monitor could be blamed for missing.

\item \textbf{A four-way recovery comparison}---no recovery, retry-only, rule-based (canonical
operators), and model-based (learned operator selection)---scored on the same episodes with
auto-verified outcomes, cluster-bootstrap confidence intervals over task templates, and a fixed
price sheet so that cost is part of the result.

\item \textbf{Release} of the dataset, deterministic harness, baseline detector and evaluator,
optional recovery ranker, reference agents, fixed configs, and a leaderboard protocol with pinned
versions (Sec.~\ref{sec:release}).
\end{enumerate}

\paragraph{What we find.} (i) Detection is easier than repair. A 13-rule monitor over
heterogeneous signals reaches AUC \AucRules{}, but even with oracle-quality detection the rule-based
policy repairs only \RecRules\% of harmful failures. (ii) Repair is class-dependent in a way that
inverts intuition: mechanically ``hard'' classes such as contradictory evidence are repaired at
96\%, while \emph{planning error} (41\%) and \emph{objective drift} (52\%)---the classes that
require knowing what the task was supposed to be---are where recovery collapses. (iii) Timing
carries real signal: episodes detected before the fault propagates are recovered \EarlyRec\% of the
time, versus \LateRec\% for late detection. (iv) A learned operator ranker trained off-policy on
logged decisions reaches \RecModel\%, i.e.\ it does not beat the hand-written canonical mapping; we
report this negative result rather than tuning it away. (v) Multi-agent episodes recover less often
than single-agent episodes on the same fault set (\MultiRecovery\% vs \SingleRecovery\%), and
coordination faults are recoverable only when the message bus is re-synchronised---a repair
operator that has no single-agent analogue.

\paragraph{What this is not.} \arb{} is entirely simulated. It does not measure production
reliability, and relative ordering under injected faults need not transfer to the messy failures of
real deployments. The reference agents are scripted workflow executors with a known standard
operating procedure and no oracle access; one experiment replaces the recovery \emph{controller}
with a local 7B language model (Sec.~\ref{sec:results}), but no frontier model plans the tasks, and
no frontier model is evaluated here. We state what this buys (exact control, exact verification,
exact replay) and what it costs (external validity) in Sec.~\ref{sec:limits}.
